% sec1.tex — c01｜原文 p1–5（并含 p6 上半的 §1.4–§1.5，即 §1 全文）｜§1 引言
% 公式号 (1)–(5)；编号环境 1.1–1.12
% 块边界：c02 从 p6 的「2. The homotopy type of quantum mechanics」标题开始，不重复本块 §1.4–§1.5。

% [原文 p.1]
20 世纪 80 年代后期，Witten 利用理论物理的方法构造了一个取值于整模形式的椭圆亏格 [Wit88a]。他的论证暗示，椭圆上同调与量子场论之间存在深得多的联系。沿着这些想法，Segal 把 2 维量子场论精确地定义为配边范畴的一个射影表示 [Seg04]。同一时期，Hopkins 及其合作者发展了椭圆上同调这一课题，其顶点是拓扑模形式（TMF）的构造，以及把 Witten 亏格提升为环谱之映射的 string 定向 [Hop02, Lur09b]。经过跨越数十年的种种洞见与努力，一幅把 TMF 与场论的几何联系起来的丰富图景正在浮现。在 Stolz 与 Teichner 的框架中，基本猜想可概述如下 [ST11]。

\begin{conjecture}[1.1]
对每个 $n\in\Z$ 与光滑流形 $X$，存在虚线箭头
\begin{equation}
\begin{tikzpicture}[baseline=(current bounding box.center), >=stealth]
  \node (A) at (0,3.6) {$\left\{\begin{array}{c}\text{一族 $n$ 维 string 流形}\\[2pt] \text{$M\to X$}\end{array}\right\}$};
  \node (C) at (8.4,1.8) {$\TMF^{-n}(X)$};
  \node (B) at (0,0) {$\left\{\begin{array}{c}\text{全扩张的、次数为 $-n$ 的}\\[2pt] \text{2 维超对称场论}\\[2pt] \text{（定义在 $X$ 上）}\end{array}\right\}$};
  \draw[->] (A.south east) -- node[above, sloped] {拓扑指标} (C.north west);
  \draw[->, dashed] (A.south) -- node[left] {量子化} (B.north);
  \draw[->, dashed] (B.east) -- node[below right=1pt, pos=0.42] {上闭链${}_{-n}$} (C.south west);
\end{tikzpicture}
\tag{1}
\end{equation}
使三角形交换；其中“拓扑指标”是 string 定向，“量子化”是对 $M\to X$ 的纤维上的超对称 $\sigma$ 模型作量子化，“上闭链$_n$”把 $X$ 上次数为 $n$ 的场论送到 $\TMF^n(X)$ 中的一个类。当 $n$ 为负时，空流形是唯一的 string 流形，此时 (1) 的内容是存在一个上闭链映射。
\end{conjecture}
% [Tutor Note: 式 (1) 中“拓扑指标”那条箭头在原文画作实线，另两条为虚线，而正文统称 "dashed arrows"；照译。实线表示 string 定向这一已知映射，虚线才是待构造的部分。]
% [Tutor Note: 式 (1) 的箭头标签是 $\mathrm{cocycle}_{-n}$（下标带负号），而紧邻的正文写作 cocycle$_n$；二者不矛盾：三角形针对 degree $-n$ 的场论，故下标随度数变号。]

简言之，猜想 1.1 提出 TMF 的一种几何构造，其上闭链是 2 维量子场论。它还提出一个 TMF-指标定理：当 $n\ge 0$ 时，虚线箭头的复合定义一个解析指标，而三角形 (1) 把它与 Ando–Hopkins–Rezk 构造的拓扑指标等同起来 [AHR10]。Stolz 与 Teichner 已经解释了这一被提议的 TMF-指标定理如何推广 Atiyah 与 Singer 的指标定理 [ST04]；见下文猜想 3.1。
% [Tutor Note: 原文 p1 底部另有一条不编号的注（作者单位与 NSF 资助 DMS-2205835、DMS-2340239），译本把它移入封面，以保持正文脚注编号与原文一致（本块即原文脚注 1、2）。]

% [原文 p.2]
\subsection{猜想为何有趣}
拓扑 K 理论在 20 世纪数学中的核心地位，来自它与诸多学科相互交织的方式。在几何中，向量丛决定 K 理论类。在分析中，Fredholm 算子空间表示 K 理论。在代数中，紧 Lie 群的表示决定（等变的）K 理论类。许多重要结果都依赖于这些相互联系，其中 Atiyah–Singer 指标定理是沟通分析、几何与拓扑的登峰造极之例。在代数拓扑内部，TMF 比 K 理论深刻得多，例如 $\pi_*\TMF$ 探测到球面稳定同伦群中多得多的类。然而，几何描述的缺失在很大程度上把 TMF 的触角限制在同伦理论的领域内。解决猜想 1.1 将改变这一局面。

与来自代数拓扑的这一动机互补，理论物理当前的一个主题是理解各种量子场论的空间，包括 string 理论中的景观问题 [Sus07, Dou10] 与物质的拓扑相 [Kit09a, FM13]。猜想 1.1 断言，其中一个这样的量子场论空间具有非常有趣的拓扑。所讨论的量子理论在 2 维具有最小超对称；这给可观测量施加的约束相对贫乏，使得标准计算变得困难。在假定猜想 1.1 成立的前提下所作的物理推理，已经产生了一连串令人惊讶的预言，其中许多又与数学相联系，例如参见 [GJF23, GPPV21, Tac21, TY23b, TY23a, TYY23, LP23]。举两个具体例子：TMF 中丰富的挠元指示出量子场论的新形变不变量，它们细化椭圆亏格 [GJFW21]；TMF 的 576 周期性则预言了手征自由费米子中一个令人惊讶的周期性 [DH11]。总而言之，人们期待解决猜想 1.1 将为研究一般的量子场论空间提供重要的数学工具，并特别地给出对 2 维场论与 string 理论结构的深刻洞察 [Wit88a]。

人们还期待 (1) 在同伦理论与理论物理之外的应用。下面强调三个方面。

在微分几何中，Stolz 猜想 Witten 亏格是正 Ricci 曲率度量存在的障碍 [Sto96]。正标量曲率的类似障碍，即 $\hat{A}$-亏格障碍，使用的是 Dirac 算子的几何 [Lic63]。与此相关，Hopkins 勾勒了如何通过 signature 算子的一个 TMF 推广来探测自由环路空间上的异种光滑结构 [DFHH14, 第 10 章]，这一推广由 Witten 提议 [Wit88a]。若不能牢牢把握 (1) 中场论的几何，这类环路空间指标理论仍将遥不可及。

在代数中，格顶点算子代数可以编码散在群，最著名的是通过 Borcherd's 对魔群月光猜想的证明 [Bor92]。
% [Tutor Note: 原文 p2 写作 "Borcherd's proof"，应为 "Borcherds' proof"（Richard Borcherds，1992 年证明 monstrous moonshine）；按“笔误照排”处理，已登记 ERRATA。]
格的 $\theta$ 函数也决定一个拓扑模形式 [Hop02, §5]。人们期待 (1) 中的上闭链映射能澄清围绕魔群、2 维共形场论与椭圆上同调的这些巧合及其他巧合 [Mor09, JF20, Lin22, AKL23]。

在表示论中，Grojnowski 描述了由等变上同调与等变 K 理论构成的标准例子的椭圆推广 [Gro07]。用复解析等变椭圆上同调（粗略地说，等变的 $\TMF\otimes\C$）推进这些想法，已经引出一些极其丰富的对象 [GKV95, ZZ15, YZ17, FRV17, RTV17, AO21, RSVZ22, GR13, KRW19, RW20, RW22]。
% [Tutor Note: 原文 p2 写 "has lead to"，应为 "has led to"；中文无形态变化，照译，已登记 ERRATA。]
% [Tutor Note: “等变的 $\TMF\otimes\C$” 指复化（complexification），而非给 TMF 增加新的张量结构；原文以 “roughly” 限定。]
在这些构造中——与 K 理论和普通上同调形成对照——相应的等变椭圆上同调类并没有被很好理解的几何描述。相应地，这些“椭圆主题”的代数的真正起源仍有些晦暗不明。对 (1) 的更好理解似乎必然会为这些深刻的对象提供进一步的洞见，例如通过整性细化。

\begin{remark}[1.2]
猜想 1.1 有一个重要的加强版本，它断言：$X$ 上次数为 $n$ 的场论的 concordance 类与 $\TMF^n(X)$ 一一对应。换言之，TMF 是 2 维超对称场论的完全形变不变量。这严格地难于猜想 1.1，并引入笔者希望在本综述中回避的技术性细节。简短的讨论参见下文注 3.21。
\end{remark}
% [Tutor Note: concordance 是场论之间的一种等价关系（沿超区间族的同伦），比猜想 1.1 的“存在上闭链”强得多；本译保留原文 concordance 并附中文“协调类”。]

% [原文 p.3]
\subsection{超对称量子场论的框架}
数学量子场论中一个反复出现的主题是算子代数分析与时空几何之间的相互作用。例如，Jones 的贡献——经 Witten 进一步放大——利用量子场论在纽结的几何与 von Neumann 代数的分析之间锻造出深刻的联系 [Jon87, Wit89, Ati90]。人们期待猜想 1.1 将沿着这条道路继续。Atiyah [Ati89] 与 Segal [Seg04] 提供了一种考察几何与分析之交界的富有影响力的语言：一个 $d$ 维量子场论是从 $d$ 维配边范畴到拓扑向量空间的对称幺半函子，
\begin{equation}
(\Bord_d,\bigsqcup)\longrightarrow(\Vect,\otimes).
\tag{2}
\end{equation}
几何位于 (2) 的源中，泛函分析位于其目标中。例如在 2 维，源可以选取得以捕捉 Riemann 面的模空间（如 Segal 关于共形场论的原始工作），目标可以选取得以编码核 Fréchet 空间在射影张量积下的分析。
% [Tutor Note: 式 (2) 的源范畴用不交并 $\bigsqcup$ 作幺半结构、目标用张量积 $\otimes$；这正是“配边→向量空间”的函子成为对称幺半函子的条件。]

除 (2) 之外，如今关于数学量子场论有许多流派，各自为 (1) 带来自己的洞见，例如顶点代数 [GMS00, HK04, GL22]、因子化代数 [Cos10, GG14] 与共形网 [Was98, BDH15]。归根结底，这些语言中每一种的发展，以及比较它们的“词典”的发展，都将加快猜想 1.1 以及数学量子场论中其他未决问题的进展。为了具体起见，下文将遵循 Stolz 与 Teichner [ST11]，通过 (2) 的视角考察猜想 1.1。

首先，对一个刚性几何 $\mathcal G$ 与一个光滑叠 $\cX$，存在 $\cX$ 上带 $\mathcal G$-结构的配边的对称幺半范畴 $\Bord_{\mathcal G}(\cX)$ [ST11, Definitions 2.21, 2.46, 2.48, and 4.4]。配边是定义量子场论所在的时空或世界面，叠 $\cX$ 则贡献背景场。在猜想 1.1 的情形中，背景场就是到光滑流形 $X$ 的映射。但允许更一般的背景场（例如带联络的 $G$-丛，即背景规范场）将是重要的。

下面两类例子可以让人对范畴 $\Bord_{\mathcal G}(\cX)$ 有所体会。

\begin{example}[1.3（欧几里得配边范畴）]
一个原型性的刚性几何是 $\R^d$ 上的平坦 spin 几何，其特征是等距群 $\R^d\rtimes\mathrm{Spin}(d)$ 通过平移与保向旋转（经由二重覆盖 $\mathrm{Spin}(d)\to\mathrm{SO}(d)$）作用在 $\R^n$ 上。此时记相应的 $\cX$ 上 $d$ 维欧几里得配边范畴为 $\EBord_d(\cX)$。$\EBord_d(\cX)$ 中的态射（记作 $\Phi\colon\Sigma\to\cX$）是带平坦度量与测地边界的紧 spin $d$-流形 $\Sigma$，连同映射 $\Phi\colon\Sigma\to\cX$。配边的源与目标来自把 $\Phi$ 限制到边界的一个划分 $\partial\Sigma=\partial_{\mathrm{in}}\Sigma\sqcup\partial_{\mathrm{out}}\Sigma$ 上，分为入射（源）与出射（目标）两部分。\footnote{为简洁起见，定义中省略了几个重要的技术细节。例如，为使 $\EBord_d$ 中的复合有良好定义，要求 $\partial_{\mathrm{in}}\Sigma$ 与 $\partial_{\mathrm{out}}\Sigma$ 具有领（collar）。细节参见 [ST11]。}
\end{example}
% [Tutor Note: 例 1.3 中刚性几何定义在 $\R^d$ 上，而等距群却写作作用在 $\R^n$ 上；原文如此（沿用 ST11 的记法，$n$ 记环境维数），此处应读作 $n=d$，照排不改。]

\begin{example}[1.4（超欧几里得配边范畴）]
把例 1.3 加以推广，超欧几里得刚性几何决定配边范畴 $\EBord_{d|\delta}(\cX)$，它依赖于一对自然数 $d|\delta$（即一个超维数）以及附加数据 [ST11, §4.2]。粗略地说，$\EBord_{d|\delta}(\cX)$ 中的态射仍是带平坦度量与测地边界的紧 spin $d$-流形 $\Sigma$，连同映射 $\Phi\colon\Sigma\to\cX$。新特征是：$\delta$ 维复旋量丛 $S\to\Sigma$ 被视为“奇的”，而超对称是 $S\to\Sigma$ 的自同构，它们把 $\Sigma$ 中的偶方向与 $S$ 中的奇方向混合起来。粗略到一阶近似，$\EBord_{d|\delta}(\cX)$ 就是 $\EBord_d(\cX)$ 再加上这些奇自同构。要把这一点说精确，需要超流形的语言 [Fre99, DM99]。
\end{example}

由此得到我们对超对称场论的第一个概念（后文还会展开）。

% [原文 p.4]
\begin{definition}[{1.5（[ST11, Definition 5.2] 之略述）}]
以 $\Alg$ 记 $\Z/2$-分次拓扑代数、双模与双模映射构成的 2-范畴，其对称幺半结构来自 $\Z/2$-分次张量积。光滑叠 $\cX$ 上的扭 $d|\delta$ 维超对称欧几里得场论是一个自然变换 $E$
\begin{equation}
\begin{tikzpicture}[baseline=(current bounding box.center), >=stealth]
  \node (A) at (0,0) {$\EBord_{d|\delta}(\cX)$};
  \node (C) at (7.6,0) {$\Alg$};
  \draw[->] (A) to[bend left=30] node[above] {$\mathbf 1$} (C);
  \draw[->] (A) to[bend right=30] node[below] {$\alpha$} (C);
  \node at (3.8,0) {$E\ \Downarrow$};
\end{tikzpicture}
\tag{3}
\end{equation}
介于常值函子 $\mathbf 1$（取值于 $\Alg$ 的幺半单位）与称为扭的函子 $\alpha$ 之间。$\cX$ 上的扭场论是一个群胚的对象。
\end{definition}
% [Tutor Note: 定义 1.5 中的 $\mathbf 1$ 是取值为 $\Alg$ 幺半单位的常值函子，不是数字 1；式 (3) 中 $E$ 是自然变换 $\mathbf 1\Rightarrow\alpha$，2-胞的上下两条 1-态射分别是 $\mathbf 1$ 与 $\alpha$。]

对定义的若干结构方面作些说明。第一，场论对 $\cX$ 是自然的：对映射 $\cX\to\cX'$，显然的函子 $\EBord(\cX)\to\EBord(\cX')$ 诱导出从 $\cX'$ 上的扭场论到 $\cX$ 上的扭场论的拉回。第二，扭与扭场论有乘法：$\Alg$ 上的幺半结构给出
\[
(E\colon\mathbf 1\Rightarrow\alpha,\ E'\colon\mathbf 1\Rightarrow\alpha')\ \overset{\otimes}{\longmapsto}\ (E\otimes E'\colon\mathbf 1\simeq\mathbf 1\otimes\mathbf 1\Rightarrow\alpha\otimes\alpha').
\]
扭 $\alpha$ 称为可逆的，若存在扭 $\alpha'$ 以及与平凡扭的等价 $\alpha\otimes\alpha'\simeq\mathbf 1$。下文考虑的所有扭都是可逆的。

可逆扭是刻画反常量子场论的一种方式 [Fre14]。粗略地说，射影对称性决定反常，而扭场论是配边范畴的射影表示。定义 1.5 的各种变体已在文献中出现（常常是独立出现的），例如 Freed 与 Teleman 的相对场论 [FT14]。Stolz 与 Teichner 所选的术语，来自反常量子场论与扭上同调类之间一个发展中的联系 [Sto19]。

\begin{remark}[1.6]
提请注意：定义 1.5 中的扭与 [Wit91] 中的拓扑扭无关，例如镜像对称中的 A-扭与 B-扭 [Wit88b]。
\end{remark}

\begin{example}[1.7（非扭场论）]
以 $\Vect$ 记 $\C$ 上 $\Z/2$-分次（拓扑）向量空间带张量积构成的对称幺半范畴。非扭场论是对称幺半函子 $\EBord_{d|\delta}(\cX)\to\Vect$，它给出平凡扭 $\alpha=\mathbf 1$ 下的自然变换 (3) [ST11, Lemma 5.7]。

扭的数据的一部分，是对每个对象 $Y\in\EBord_{d|\delta}(\cX)$ 的一个代数 $\alpha(Y)$；扭场论的数据的一部分，是一个左 $\alpha(Y)$-模 $E(Y)$。这样，扭场论可以看作取值于 $\alpha(Y)$-模（而非向量空间）的非扭场论；参见 [ST11, page 52]。
\end{example}

\begin{example}[1.8（次数扭）]
在定义 1.5 中取 $d|\delta=1|1$ 与 $2|1$，则对每个 $n\in\Z$ 存在一个次数扭 [ST11, §5.3]（另见下文定义 3.18 与 4.15）
\begin{equation}
\mathrm{EBord}_{d|1}(\mathrm{pt})\xrightarrow{\ \alpha_n\ }\Alg,
\qquad \alpha_n\otimes\alpha_m\simeq\alpha_{n+m},
\qquad n,m\in\Z,\ d\in\{1,2\}.
\tag{4}
\end{equation}
当 $d=1$ 时，第 $n$ 个次数扭由第 $n$ 个 Clifford 代数 $\mathrm{Cl}_n$ 决定，见式 (37)。当 $d=2$ 时，它由 $n$ 个（手征）自由费米子决定，其中包括无穷维 Clifford 代数 $\mathrm{Cl}(L\R^n)$ 的数据，其中 $L\R^n$ 是 $\R^n$ 的自由环路空间，见式 (54)。
\end{example}
% [Tutor Note: 次数扭只对 $d|1$（$d=1,2$）定义；$d=2$ 时由 $n$ 个手征自由费米子给出，涉及 $L\R^n$ 上的无穷维 Clifford 代数，对应式 (54)。]

为准确捕捉所期望的例子，我们给扭场论增添一种附加的结构与性质，称为反射正性。简言之，这是关于扭与扭场论的 $\Z/2$-等变数据，相对于配边上反转定向的对合，以及由 $\C$ 上复共轭给出的 $\Alg$ 上的对合。相对于这些等变数据，理论中的 Hermite 配对被要求是内积。综述性讨论见 [FH21, §3] 或 [KS21, page 18]。在定义 1.5 的语境下，反射正性的完整定义仍在发展之中；一个可用的定义出现在 [BE23c, §C.3] 中，我们在下文定义 2.13 中勾勒相关特征。
% [Tutor Note: 反射正性的等变数据含两处对合：配边上的定向反转、$\Alg$ 上由复共轭诱导的对合；二者作用后 Hermite 配对须为正定内积。]

% [原文 p.5]
\begin{definition}[1.9]
以 $\QFT^{\alpha}_{d|\delta}(\cX)$ 记如下群胚：其对象是 $\cX$ 上反射正的 $\alpha$-扭 $d|\delta$ 维欧几里得场论，其态射是与反射结构相容的扭场论的同构。
\end{definition}

\begin{definition}[1.10]
对 $n\in\Z$ 与光滑叠 $\cX$，以
\[
\QFT^{n}_{d|1}(\cX):=\QFT^{\alpha_n}_{d|1}(\cX)
\]
记由 (4) 沿映射 $\cX\to\mathrm{pt}$ 的拉回所扭的场论构成的范畴。
\end{definition}

特别地取光滑流形 $\cX=X$ 时，范畴 $\QFT^{n}_{d|1}(X)$ 是猜想 1.1 中所期望的对象的最著名近似。
% [Tutor Note: 原文此处写 "Specializing to a smooth manifold X = X"，其中前一个 X 实为叠 $\cX$；原文记号混用，照译作 $\cX=X$。]
% [Tutor Note: 原文 "and (at present) it unclear which one to take" 缺系动词，应为 "it is unclear"；照译，已登记 ERRATA。]
但若要对证明该猜想抱有任何希望，还需要进一步的修改：切除要求全扩张的场论，见 [ST11, Conjecture 1.17] 与 [ST04, §1]。这一向下扩张的过程需要超欧几里得配边的 2-范畴，以及把 $\Alg$ 去环（deloop）的 3-范畴。定义 1.9 并未唯一确定源范畴或目标范畴。相反，存在大量的推广，而（目前）尚不清楚该取哪一个。

\subsection{用 Lurie 的高阶范畴等变性加强猜想}
缩小可能定义之“动物园”的一条途径，是在猜想 1.1 中要求更多结构。这使人更难猜出“错误”的场论定义。一个有希望的选择是 Lurie 的高阶范畴等变椭圆上同调，在 [Lur09b, §5] 中称为 2-等变椭圆上同调。它对每个紧 Lie 群 $G$ 给出一个导出概形 $\cM_G$，对每个级别 $[\ell]\in H^4(BG;\Z)$ 给出 $\cM_G$ 上的线丛 $\cL_\ell$。$\cM_G$ 底层的普通概形是椭圆曲线上平坦 $G$-丛的模空间，$\cL_\ell$ 底层的普通线丛以非 Abel $\theta$ 函数为截面。正如导出椭圆曲线模叠上函数层的整体截面给出 TMF，$\cL_\ell$ 在 $\cM_G$ 上的整体截面决定 2-等变 TMF，记作 $\TMF^{\ell}_{G}$ [Lur19, GM23]。

紧 Lie 群 $G$ 作用在流形 $X$ 上决定叠 $\cX=X/\!/^{\nabla}G$，\footnote{栈 $X/\!/^{\nabla}G$ 分类带联络的 $G$-丛 $(P,\nabla)$ 以及 $G$-等变映射 $P\to X$。丢弃联络数据给出遗忘函子 $X/\!/^{\nabla}G\to X/\!/G$，打到通常的商栈。}从而给出猜想 1.1 的 $G$-等变加细，比较 [ST11, §1.7]。
% [Tutor Note: $X/\!/^{\nabla}G$ 中的 $\nabla$ 表示保留联络数据；忘掉联络即得通常的商叠 $X/\!/G$。这是本文“背景规范场”的几何来源。]

\begin{conjecture}[1.11]
对作用在流形 $X$ 上的紧 Lie 群 $G$ 与级别 $[\ell]\in H^4(BG;\Z)$，存在推广 (1) 的上闭链映射
\begin{equation}
\begin{tikzpicture}[baseline=(current bounding box.center), >=stealth]
  \node (A) at (0,0) {$\QFT^{\ell}_{2|1}(X/\!/^{\nabla}G)$};
  \node (C) at (8.2,0) {$\TMF^{\ell}_{G}(X)$};
  \draw[->, dashed] (A) -- node[above] {上闭链} (C);
\end{tikzpicture}
\tag{5}
\end{equation}
取值于 $X$ 的 $\ell$-扭 $G$-等变 TMF。此外，对 $G$-等变 string 流形的族，存在交换三角形 (1) 的类似物。
\end{conjecture}

用物理的语言说，范畴 $\QFT^{\ell}_{2|1}(\mathrm{pt}/\!/^{\nabla}G)$ 包含带有背景规范场与指定 't Hooft 反常的场论。由一般理论 [Fre87]，这样的反常决定亏格 1 曲面（即椭圆曲线）上 $G$-丛的模空间上的线丛。这些线丛是 Grojnowski 扭复解析等变椭圆上同调的核心 [Gro07, GKV95]。初步的研究表明，$\mathrm{pt}/\!/^{\nabla}G$ 上的扭欧几里得场论决定这类线丛的截面 [BET21, BET19]。因此，在与 Lurie 框架比较时，人们期待 $\QFT^{1}(-/\!/^{\nabla}G)$ 编码导出概形 $\cM_G$，而 $\QFT^{\ell}(-/\!/^{\nabla}G)$ 编码 $\theta$-线丛的截面 $\Gamma(\cM_G,\cL_\ell)$。

除了限制全扩张场论的可能定义之外，猜想 1.11 还照亮了构造场论与 TMF 之间比较映射的一条路径：Lurie 勾勒了 2-等变椭圆上同调的一个泛性质 [Lur09b, §5.5]。猜想 1.11 为场论提出一种 2-等变结构，于是 Lurie 的泛性质将引出所期望的比较映射。这样，高阶等变性可以看作开启猜想 1.1 的关键要素。

\begin{remark}[1.12]
等变椭圆上同调与 2 维规范理论之间的联系并不新鲜。最早的想法或许属于 Grojnowski [Gro07] 与 Segal [Seg88, Theorem 5.3] [Seg07, Definition 2.4]，不过其他人也独立地发现了类似的平行关系。例如，环路空间上 Dirac 算子的 Witten 刚性定理可以用 $S^1$-等变椭圆上同调来理解 [BT89, Ros01, AB02, Ros03]；而对有限群，人们早已清楚 $G$-等变椭圆上同调与 string 理论中的扭扇区（twisted sectors）相关 [DHVW85, Seg88, Dev96, Dev98, HKR00, Mor06]。

% [原文 p.6]
环路群及其正能表示也同时出现在 2 维场论（例如 WZW 模型）与等变椭圆上同调中 [Liu96, Kit09b, FHT11b, Kit14, Gan14, Lue22]。等变性还引向自然的椭圆幂运算，它们与 string 理论的几何有明显联系 [DMVV97, And00, Gan07, Gan06, Gan09, Hua18, HS20, BBES22]。
\end{remark}

\subsection{免责声明与概览}
本章有两个主要目标：(i) 为猜想 1.1 提供动机与直观；(ii) 描述其验证中仍然存在的中心困难。这两个目标都要求理解来自理论物理的例子，例如 $\sigma$ 模型与自由费米子。不幸的是，把这些（部分构造出来的）例子塞进量子场论的数学框架，会引入许多技术性细节，它们往往掩盖更大的图景。为此，下文的每一节都从某种特设的（ad hoc）量子力学系统与量子场论定义开始，这些定义明确捕捉了物理文献中出现的分析数据。特别地，对超对称量子力学，一个理论的相关数据是完全无歧义的，例子也已被很好地理解。把这些特设定义与代数拓扑中的对象联系起来之后，我们把分析数据安置在范式 (3) 的语境中，作为配边范畴的表示的一部分。这个配边定义指向 2 维量子场论中的一些新结构，它们难以仅凭超对称量子力学的朴素推广而窥见。为强调这一数学框架仍在发展之中，我们把带有证明梗概的陈述标记为假设（Hypothesis）；所引文献表明，已有密切相关的陈述被证明，但带有我们在此不讨论的技术性告诫。我们期待任何合理的定义都应当使所述的假设成立，因此本论述也意在提供具体的路标，以检验所提出的（全扩张）定义是否走在正确的轨道上。

下面简要概述本章。第 2 节先从分析的视角介绍量子力学，再把其基本结构翻译到配边框架中。第 3 节对超对称量子力学重复这一做法，并勾勒猜想 1.1 的 K 理论变体的一个版本。第 4 节从把量子场论看作量子力学之推广的分析视角开始，再把这些结构翻译到配边框架中。这一框架引出关于猜想 1.1 最终内容的各种洞见。
% [Tutor Note: 原文 "translating these structure into the bordism framework" 中 these structure 应为 this structure；中文不受影响，已登记 ERRATA。]

全文中，$\cH$ 表示一个（固定的）无穷维可分复 Hilbert 空间。$\Z/2$-分次对象的所有张量积都是 $\Z/2$-分次张量积，其对称同构编码 Koszul 符号。
% [Tutor Note: $\Z/2$-分次张量积的对称同构含 Koszul 符号，即 $a\otimes b=(-1)^{|a||b|}b\otimes a$；这是全文一切“超”结构的统一约定。]

\subsection{致谢}
感谢 Emily Cliff、André Henriques、Theo Johnson-Freyd、Matthias Ludewig、Laura Murray、Nat Stapleton 与 Arnav Tripathy 的许多富有启发的交谈，这些想法在其中成形。也感谢 Theo Johnson-Freyd 与 Laura Murray 对先前草稿的反馈。Stephan Stolz 与 Peter Teichner 的影响怎么强调都不过分——感谢他们的诸多洞见与慷慨分享想法。
